%==================================================================
% ABSTRACT (Bahasa Inggris)
%==================================================================
\phantomsection
\addcontentsline{toc}{chapter}{ABSTRACT}
\begin{center}
    \textbf{\judulen}\\[0.5cm]
    by\\
    \penulis\\
    NIM: \nim\\[1.5em]
    \textbf{ABSTRACT}\\[0.5cm]
\end{center}

\textit{The academic information system is a service used by all students, yet
complaints regarding its ease of use recur in every study plan registration
period. Its administrators lack a measurable basis for determining which aspects
should be improved first. This study aims to measure the service quality of the
academic information system, analyse the gap between importance and performance
levels, and determine improvement priorities. Measurement was conducted using
the WebQual 4.0 instrument comprising 22 items across three dimensions:
usability, information quality, and service interaction quality. Data were
collected through a survey of 96 active students and analysed descriptively,
followed by Importance-Performance Analysis to map each item into four priority
quadrants. The results show that the usability dimension obtained the lowest
performance score, while the largest gap occurred in the navigational ease item.
The quadrant mapping placed five items in the top priority quadrant, all related
to navigation and clarity of the registration flow. This study concludes that
improvement efforts should be directed towards interface information
architecture rather than the addition of new features, and produces a list of
improvement priorities that administrators can act upon directly.}\\[0.6cm]

\noindent \textbf{Keywords:} \textit{\keywords}
\clearpage
